\documentclass[10pt,a4paper]{amsart}
\usepackage[margin=19mm]{geometry}
\usepackage{amsmath,amssymb,mathtools}
\usepackage[scr=rsfs,cal=euler]{mathalfa}
\usepackage{xfrac}
\usepackage[unicode,hidelinks,bookmarks=false]{hyperref}
\newcommand{\ie}{i.e.\ }
\newcommand{\cf}{cf.\ }
\newcommand{\resp}{respectively\ }
\newcommand{\Subsection}{\subsection}
\providecommand{\nobreakdash}{\nobreak\hbox{-}}
\setlength{\emergencystretch}{2em}
\multlinegap0pt
\usepackage{amssymb}
\usepackage{enumerate}
\usepackage{tcolorbox}
\usepackage{mdframed}
\usepackage{mathtools}
\newtheorem{theorem}{Theorem}
\title{Superelliptic Affine Lie algebras and orthogonal polynomials II}
\author{Felipe Albino dos Santos, Mikhail Neklyudov, Vyacheslav Futorny}
\date{Source: arXiv:2603.29082v1 (CC BY 4.0)}
\begin{document}
\maketitle

\noindent\emph{Source and licence.} Superelliptic Affine Lie algebras and orthogonal polynomials II. Version arXiv:2603.29082v1 (CC BY 4.0), distributed by arXiv under the Creative Commons Attribution 4.0 International licence, http://creativecommons.org/licenses/by/4.0/, as recorded in the arXiv licence metadata for this exact version and shown on the abstract page https://arxiv.org/abs/2603.29082v1. Full licence notice: \texttt{licenses/arxiv-2603.29082-CC-BY-4.0.txt}.\par\smallskip

\noindent\textit{Editorial scope.} Counted unit: the article's theorem thm:orthogonality (source0.tex lines 834--839). The article's complete original proof of it is delivered with it (source0.tex lines 841--863, one \texttt{proof} environment of 23 lines). No mathematics is written, repaired or re-derived: the statement and the proof are the article's own lines, in the article's own notation, and no retained line is substituted or re-flowed.  Retained uncounted as the source-local context the proof needs: lines 773--773; lines 828--831.  Nothing was omitted from any retained span, so the package carries no omission record.  No retained line contains a citation, so the wrapper carries no bibliography and no bibliography entry.  No retained line contains a figure environment, an image-inclusion command or drawing code, so no image is delivered and the package declares no asset dependency.  Editorial formatting only, in addition: an AMS \texttt{amsart} preamble that restates the article's own declarations and macros used by the retained lines, namely the environments \texttt{theorem} on separate counters, together with the standard packages those lines need.  The retained environments print in this document as Theorem 1 in order of appearance, whereas the article prints the same items with the numbering of its own full document.  The complete native source of the article as distributed in the arXiv e-print for this exact version is bundled unchanged as \texttt{sources/197491/source0.tex}.
\section{Associated ultraspherical polynomials}\label{subsec:assoc-ultra}

\begin{equation}\label{eq:assoc-ultra-rec-orth}
2c\,(n+\nu+c_0)\,q_n(c)=(n+c_0)\,q_{n-1}(c)+(n+2\nu+c_0)\,q_{n+1}(c),
\qquad n\ge0,
\end{equation}

\begin{theorem}\label{thm:orthogonality}
Assume $r\ge2$ and $m\ge2$. Then there exists a positive Borel measure $\mu$ on $[-1,1]$
such that $\{q_n(c)\}$ is an orthogonal polynomial system in $L^2([-1,1],d\mu)$.
Consequently, the superelliptic families in cases~1 and~2 yield orthogonal polynomial
systems after parity restriction and the reindexing of Section~\ref{subsec:assoc-ultra}.
\end{theorem}

\begin{proof}
Rewrite \eqref{eq:assoc-ultra-rec-orth} in the form
\begin{equation}\label{eq:xqn}
c\,q_n(c)=A_n\,q_{n+1}(c)+B_n\,q_{n-1}(c),
\qquad
A_n=\frac{n+2\nu+c_0}{2(n+\nu+c_0)},\quad
B_n=\frac{n+c_0}{2(n+\nu+c_0)}.
\end{equation}
Since $\nu>0$ and $c_0>0$, we have $A_n>0$ for all $n\ge0$ and $B_n>0$ for all $n\ge1$.
Define a renormalized sequence $\{p_n\}$ by $p_n=\gamma_n^{-1}q_n$, where $\gamma_0=1$ and
$\gamma_{n+1}=A_n\gamma_n$. Multiplying \eqref{eq:xqn} by $\gamma_n^{-1}$ yields the
Favard form
\begin{equation}\label{eq:favard}
c\,p_n(c)=p_{n+1}(c)+a_n\,p_{n-1}(c),\qquad n\ge0,
\end{equation}
with
\[
a_n=B_nA_{n-1}=\frac{(n+c_0)(n-1+2\nu+c_0)}{4(n+\nu+c_0)(n-1+\nu+c_0)}>0\quad(n\ge1).
\]
By Favard's theorem, \eqref{eq:favard} implies the existence of a positive Borel measure
$\mu$ for which $\{p_n\}$ is orthogonal. Since $p_n$ differs from $q_n$ only by a nonzero
scalar, $\{q_n\}$ is orthogonal with respect to the same measure.
\end{proof}

\end{document}
